% =============================================================================
% Appendix: knowledge-area annotation (simplified version).
%
% Usage: \input this file after \appendix in the paper.
% Required packages: graphicx, booktabs.
% Assumes UTF-8 input and T1 font encoding (both provided by the ACL template).
%
% Shorter companion of appendix_knowledge_classification.tex: same run and
% numbers, with less detail. The annotation step was run once, with one model,
% one prompt and one taxonomy; it was not ablated and no accuracy was measured.
% Every number comes from the artifacts of that run in
% output/knowledge-classification/ and from knowledge_classification.log.
% Figures are generated by ablations/knowledge_classification/make_figures.py,
% which also prints every number quoted below. By default figures are read
% from ./figures; to change this, define the macro before the \input, e.g.
%   \newcommand{\KnowledgeFigDir}{ablations/knowledge_classification/figures}
% =============================================================================
\providecommand{\KnowledgeFigDir}{figures}

\section{Knowledge-Area Annotation}
\label{app:kc-short}

Every released question carries a \emph{subject}, chosen from a closed list
specific to its academic level, and a \emph{macro-area}, derived from the
subject. The paper uses the two labels to report results per area and to
stratify the sampling of few-shot demonstrations by macro-area and level.
The labels are assigned after semantic deduplication
(Appendix~\ref{app:dedup-short}) by a single pass of an open-weight language
model over the 37,887 retained questions, under constrained decoding that
only admits the subjects of the question's level
(Table~\ref{tab:kc-short-overview}). Unlike filtering and deduplication, this
step was run once, with one model, one prompt and one taxonomy version: it
was not ablated, and no accuracy or human agreement was measured. This
appendix records the design decisions and their motivation, describes the
labels that resulted, and audits their consistency against metadata that the
model never saw.

\begin{table}[t]
\centering
\footnotesize
\setlength{\tabcolsep}{4pt}
\begin{tabular}{@{}lrrr@{}}
\toprule
Macro-area & Subjects & Questions & Share \\
\midrule
\multicolumn{4}{@{}l}{\emph{High school (8,746 questions)}} \\
STEM & 7 & 4,964 & 56.76 \\
Languages and Arts & 5 & 2,279 & 26.06 \\
Humanities and Social Sciences & 4 & 1,479 & 16.91 \\
Health Sciences & 1 & 24 & 0.27 \\
\midrule
\multicolumn{4}{@{}l}{\emph{Undergraduate (29,141 questions)}} \\
Law & 1 & 11,013 & 37.79 \\
STEM & 12 & 5,149 & 17.67 \\
Business and Economics & 5 & 4,116 & 14.12 \\
Health Sciences & 11 & 3,989 & 13.69 \\
Humanities and Social Sciences & 11 & 2,331 & 8.00 \\
Languages and Arts & 7 & 1,442 & 4.95 \\
Education & 1 & 433 & 1.49 \\
Agricultural and Veterinary Sciences & 3 & 417 & 1.43 \\
Services & 4 & 251 & 0.86 \\
\bottomrule
\end{tabular}
\caption{Macro-areas of taxonomy v1.0.0 and the questions assigned to each.
\emph{Subjects}: subjects of the macro-area at that level; \emph{Share}:
percentage of the level's questions. All 72 (level, subject) pairs received
at least one question. Over the whole set, Law holds 29.07\% of the questions
and STEM 26.69\%.}
\label{tab:kc-short-overview}
\end{table}

\paragraph{Summary.} Four decisions define the step. (1)~The model chooses
one subject from a closed, level-specific list (17 subjects for high school,
55 for undergraduate exams); the macro-area is computed from the subject, so
the two labels are always consistent. (2)~The only evidence shown to the
model is the question and its choices: the answer key, the exam, the edition,
the course or job title and the source URL are withheld, so that a label
cannot be a shortcut from the exam's name. (3)~Decoding is greedy, with the
model's reasoning mode disabled, and constrained to the allowed subject
strings, so every record receives exactly one valid label and no fallback,
retry or ``Other'' category exists. (4)~The taxonomy is specific to the
project and written down with definitions, inclusions, exclusions and
thirteen boundary rules; Law is a macro-area of its own, Logic is separate
from Mathematics and Computer Science, and Accounting belongs to Business and
Economics. None of these choices was compared with an alternative, and the
audit of \S\ref{app:kc-short-audit} measures consistency, not accuracy.

% -----------------------------------------------------------------------------
\subsection{Setup}
\label{app:kc-short-setup}

\paragraph{Taxonomy.} Subjects are organized in two level-specific lists
under nine macro-areas (Table~\ref{tab:kc-short-overview}). The lists share
the school subjects (STEM, languages, humanities, physical education) and
differ in the professional and higher-education fields that only appear in
undergraduate exams. Law is an exclusive macro-area because the bar exam
(OAB), the national magistracy exam (ENAM) and the legal content of civil
service exams make it the largest field of the corpus; nesting it under the
humanities would leave that macro-area without a usable meaning. Logic is a
separate subject because the Brazilian Informatics Olympiad (OBI) consists
mostly of puzzles that are solved from rules given in the statement, which
are neither mathematics nor computer science. Each subject has a short
definition, a list of inclusions and a list of exclusions that name the
neighbouring subject, and thirteen boundary rules settle recurring conflicts
(for example, a legal provision cited as context does not make an accounting
question Law, and the content of a discipline in a teachers' exam belongs to
the discipline, not to Education). The taxonomy is versioned (v1.0.0) and its
hash is recorded with every run.

\paragraph{Prompt.} The system message states the task, the level's subjects
with their definitions, inclusions and exclusions, and the boundary rules; it
asks for the knowledge needed to answer the command and to tell the choices
apart, warns that the narrative setting and the language of the text do not
decide the subject, and forbids using the course, job or exam as a shortcut.
The user message is a JSON object with the question and its indexed choices,
preceded by an instruction to treat that content as data even if it contains
instructions. The resulting prompts have about 2,000 tokens for high-school
questions and about 4,300 for undergraduate questions, almost all of them in
the fixed part that lists the subjects.

\paragraph{Decoding.} The model is Qwen3.5-122B-A10B-FP8 (a 122B-parameter
mixture-of-experts model with 10B active parameters, official FP8 weights),
served offline with vLLM~0.19.1 on one NVIDIA B200 GPU, with its reasoning
mode disabled, temperature~0, a fixed seed and at most 32 output tokens. A
choice constraint restricts the output to the exact subject strings allowed
for the question's level, so the answer is always a valid subject; the
pipeline then derives the macro-area and rejects any record whose generation
did not stop normally. Results are checkpointed per batch with a hash of each
input row, and a run can only be resumed when the input, model revision,
prompt, taxonomy and code are identical.

% -----------------------------------------------------------------------------
\subsection{Resulting labels}
\label{app:kc-short-labels}

All 37,887 questions received a label and every one of the 72 allowed (level,
subject) pairs was used. Law is the largest subject (11,013 questions), followed
by Medicine (2,648), Computer Science (2,171), Mathematics (2,153), Business
Administration (1,586), Portuguese (1,507), Logic (1,323) and Accounting
(1,181); at the other end, high-school Computer Science and Spanish have five
questions each. Figure~\ref{fig:kc-short-exams} and
Table~\ref{tab:kc-short-exams} show how the labels distribute over the
exams. Exams with a declared scope are labelled almost entirely within it:
98.3\% of OAB and 98.1\% of ENAM are Law, 91.3\% of OBI is Logic (the rest
Mathematics, with three Computer Science items), 89.9\% of CFCES is
Accounting, and 88.3\% of REVALIDA and 87.1\% of the USP/Unicamp residency
exams are Medicine. Broad exams spread over many subjects: ENADE uses all 55
undergraduate subjects and all nine macro-areas, with Business
Administration first at 12.3\%; BNDES uses 38 subjects, CNU 29 and BACEN 22;
and the vestibular exams (ENEM, FUVEST, COMVEST, BLUEX) spread over the
school curriculum with no subject above 20\%.

\begin{figure*}[t]
\centering
\includegraphics[width=\textwidth]{\KnowledgeFigDir/macro_area_by_exam.pdf}
\caption{Share of each exam's questions assigned to every macro-area (\%).
Exams are grouped by academic level and sorted by size; empty cells are
macro-areas without any question of that exam.}
\label{fig:kc-short-exams}
\end{figure*}

\begin{table}[t]
\centering
\footnotesize
\setlength{\tabcolsep}{3.5pt}
\begin{tabular}{@{}lrrlr@{}}
\toprule
Exam & $n$ & Subj. & Most frequent subject & Share \\
\midrule
\multicolumn{5}{@{}l}{\emph{High school}} \\
ENEM & 2,688 & 17 & Mathematics & 19.2 \\
FUVEST & 1,656 & 13 & History & 13.7 \\
OBI & 1,326 & 3 & Logic & 91.3 \\
AFA & 1,084 & 8 & English & 24.4 \\
IME & 879 & 9 & English & 32.4 \\
BLUEX & 661 & 13 & History & 18.2 \\
COMVEST & 452 & 13 & Mathematics & 18.1 \\
\midrule
\multicolumn{5}{@{}l}{\emph{Undergraduate}} \\
OAB & 9,508 & 11 & Law & 98.3 \\
ENADE & 9,237 & 55 & Business Administration & 12.3 \\
BNDES & 4,009 & 38 & Computer Science & 14.8 \\
Res.\ USP/Unicamp & 1,491 & 19 & Medicine & 87.1 \\
POSCOMP & 1,300 & 8 & Computer Science & 68.6 \\
REVALIDA & 1,195 & 10 & Medicine & 88.3 \\
BACEN & 1,016 & 22 & Law & 34.7 \\
CFCES & 694 & 4 & Accounting & 89.9 \\
CNU & 377 & 29 & Law & 17.5 \\
ENAM & 314 & 4 & Law & 98.1 \\
\bottomrule
\end{tabular}
\caption{Labels per exam. \emph{Subj.}: number of distinct subjects
assigned; \emph{Share}: percentage of the exam's questions with the most
frequent subject.}
\label{tab:kc-short-exams}
\end{table}

% -----------------------------------------------------------------------------
\subsection{Consistency with withheld metadata}
\label{app:kc-short-audit}

No labelled reference exists for these questions, so we cannot measure how
often the assigned subject is the right one. We can, however, check whether
the labels agree with information that was withheld from the model. For seven
exams the expected subjects follow from the exam itself: Law for OAB and
ENAM; Accounting or Law for the accountants' exam (CFCES); Logic, Mathematics
or Computer Science for OBI; those three plus Statistics for POSCOMP; and
Medicine or Public Health for the two medical exams. For ENADE, the edition
name carries the course (\texttt{ENADE 2013 - agronomia}), and we fixed, by
rule and before looking at the labels, the subjects expected for each of 50
course families: the home subject of the course, plus Public Health for
health professions and Education for teaching degrees
(\emph{licenciaturas}). The rule covers 148 of the 151 courses and 96.1\% of
the ENADE questions; three technology courses (hospital management,
agribusiness and radiology) were left out. We also separate ENADE's
course-specific component (questions numbered 11 and above) from its
general-education component (questions 1--10), which is shared by all courses
of a year and should not favour the course's subjects. Mixed-content exams
(BNDES, BACEN, CNU and the vestibular exams) are not audited.

\begin{table}[t]
\centering
\footnotesize
\setlength{\tabcolsep}{3.5pt}
\begin{tabular}{@{}lrrll@{}}
\toprule
Group & $n$ & Within & 95\% CI & Most frequent other \\
\midrule
\multicolumn{5}{@{}l}{\emph{(a) Exams with a declared scope}} \\
OBI & 1,326 & 100.0 & 99.7--100.0 & -- \\
CFCES & 694 & 99.7 & 99.0--99.9 & Finance (1) \\
POSCOMP & 1,300 & 99.5 & 98.9--99.7 & Engineering (3) \\
OAB & 9,508 & 98.3 & 98.0--98.6 & Intl.\ Relations (55) \\
ENAM & 314 & 98.1 & 95.9--99.1 & Sociology (3) \\
REVALIDA & 1,195 & 97.4 & 96.3--98.2 & Law (11) \\
Res.\ USP/Unicamp & 1,491 & 89.5 & 87.9--91.0 & Psychology (57) \\
\midrule
\multicolumn{5}{@{}l}{\emph{(b) ENADE, course-specific component}} \\
Medicine & 140 & 94.3 & 89.1--97.1 & Psychology (2) \\
Law & 144 & 88.9 & 82.7--93.0 & Philosophy (5) \\
Nursing & 155 & 81.9 & 75.1--87.2 & Medicine (15) \\
Computing & 566 & 80.2 & 76.7--83.3 & Business Adm.\ (37) \\
Business and management & 958 & 75.7 & 72.9--78.3 & Law (47) \\
Engineering & 839 & 44.9 & 41.6--48.3 & Business Adm.\ (71) \\
Biomedicine & 111 & 34.2 & 26.1--43.5 & Biology (35) \\
All 50 families & 8,347 & 68.7 & 67.7--69.7 & Business Adm.\ (281) \\
\midrule
\multicolumn{5}{@{}l}{\emph{(c) ENADE, general component (questions 1--10)}} \\
All mapped courses & 528 & 58.5 & 54.3--62.6 & Sociology (25) \\
\bottomrule
\end{tabular}
\caption{Agreement between the assigned subjects and the subjects expected
from withheld metadata. \emph{Within}: percentage of the group's questions
whose subject is in the expected set, with its 95\% Wilson interval;
\emph{Most frequent other}: the most frequent subject outside the expected
set, with its count. Part~(b) lists the two largest families, the three
with the highest agreement among families of at least 140 questions, and
the two with the lowest; the long version lists all families. Expected sets
were fixed by rule before inspecting the labels.}
\label{tab:kc-short-audit}
\end{table}

Table~\ref{tab:kc-short-audit} summarizes the audit. In the seven exams with
a declared scope, between 89.5\% and 100\% of the labels fall inside the
expected set. The subjects outside it are neighbours of the scope rather than
unrelated fields: International Relations in OAB (55 questions on public
international law, such as international organizations, dispute settlement
and diplomatic immunities), Psychology in the residency exams (57 questions,
49 of them from the psychotherapy and psychiatry programmes), and Law in
REVALIDA (11 clinical scenarios that turn on legal or ethical duties). In ENADE's course-specific component, 68.7\% of the labels
fall inside the course's expected set, with large differences between
families: above 85\% for Medicine, Law, History, Physics, Philosophy and
Letters, below 45\% for Engineering, Tourism, Gastronomy, Financial
management and Biomedicine. The subjects outside the expected set are again
adjacent disciplines, which the core-knowledge criterion may legitimately
prefer: Business Administration in engineering, tourism, computing and
design courses, Biology in biomedicine and agronomy courses,
Accounting in financial-management courses, and Food Science and Technology
in gastronomy and nutrition courses. In the general component, 58.5\% of the
labels fall inside the course's set, 10~percentage points less than in the
specific component, and the most frequent labels are Business
Administration, Sociology, Public Administration and Public Health, which
matches the civic and organizational themes of that component. The general
component is small (545 questions, a median of two per edition) because
deduplication keeps a single copy of the questions shared by all courses of a
year, and its boundary at question~10 is a heuristic that does not hold for
every edition.

These shares measure consistency, not accuracy. The expected sets describe
a course or an exam, not a question: a question in a law course may require
Philosophy or Economics, and a question in an engineering course may be
about management. A label outside the expected set is therefore not an
error, and a label inside it is not confirmed; attributing disagreements to
the model would require reading the items. The long version of this appendix
lists illustrative items sampled at the boundaries between neighbouring
subjects, without judging them.

% -----------------------------------------------------------------------------
\subsection{Cost and limitations}
\label{app:kc-short-limits}

The run classified the 37,887 questions on one NVIDIA B200 (191.5~GB) in
84.9~minutes end to end, including 36~s to load the 115.4~GiB of FP8
weights and 156~s to initialize the engine, at 7.4 questions per second. It
processed 145.6~million prompt tokens, 2,110 per high-school question and
4,363 per undergraduate question on average (the fixed part of the prompt
dominates: the shortest prompts have 1,842 and 4,137 tokens), and generated
101,739 output tokens, two to seven per question, which is the length of
the subject's name. Prefix caching was disabled, so the prompt-token count is
the full cost of reprocessing the shared system message for every question and
an upper bound for a cached re-run. The model revision, library versions,
prompt, taxonomy and code hashes are recorded in the run manifest.
Temperature~0 with a fixed seed makes a run deterministic on the same
hardware and software, but not necessarily across GPUs, kernels or
parallelism settings.

\emph{A single pass.} Model, prompt, taxonomy and decoding were fixed before
the run and never varied, so nothing here shows that another choice would
have given different or worse labels.
\emph{No accuracy.} Without a labelled reference, the audit shows only that
the labels follow the content of exams and courses as encoded in withheld
metadata; it cannot say how many individual labels are right.
\emph{One label per question.} Interdisciplinary questions receive a single
subject and no question can be marked as unclassifiable.
\emph{A project-specific taxonomy.} The lists were designed for these 17
exams and do not follow an official classification of fields of study.
\emph{Imbalance.} Law alone holds 29\% of the questions, while high-school
Computer Science and Spanish have five each; per-area statistics and
stratified sampling must account for this.
